\section{Open proof obligations and recommended research order}
\label{sec:open-problems}

The hybrid is convergent, and its tail can be placed on the desired accelerated
work scale.  The unresolved issue is the locality of the accelerated burn-in
and, for RPPR, the local cost of a momentum-free cleanup from an arbitrary
accelerated warm start.  The following research order minimizes logical
dependence and empirical risk.

\subsection{Priority 1: establish or refute early AESP locality}

\begin{problem}[Early-locality lemma]
\label{prob:early-locality}
For the AESP run truncated at the objective-gap handoff stage \(J\), prove a
bound of the form
\begin{equation}
    \max_{t\leq J}\Lambda_t
    \leq\frac{K}{\eps},
    \label{eq:open-early-locality}
\end{equation}
with graph-independent or structurally interpretable \(K\); alternatively,
construct a graph family showing that no such uniform bound is possible for
the current FIFO AESP implementation.
\end{problem}

A useful intermediate target is a boundary-sensitive summation rather than a
uniform maximum:
\begin{equation}
    \sum_{t=1}^{J}\Work_t^{\rm in}
    \leq
    \softO\!\left(
       \frac{1}{\sqrt\alpha\eps}
       +\mathsf{BoundaryOverhead}
    \right).
    \label{eq:boundary-sensitive-target}
\end{equation}
This would parallel the core-plus-spurious-work decomposition of
\citet{fountoulakis2026complexity} and may be more realistic than requiring
all active sets to remain uniformly small.

\subsection{Priority 2: adapt KKT slack to changing shifted systems}

The over-regularized RPPR proof obtains a uniform inactive-coordinate margin
by comparing two regularization levels.  An AESP analogue would need to
address three new effects:
\begin{enumerate}[leftmargin=2em]
\item the linear term changes with the outer center \(y^{(t-1)}\);
\item the inner objective is shifted by \(\eta I\);
\item the relevant activation threshold \(\eps_t\) changes with the Catalyst
      accuracy schedule.
\end{enumerate}
A successful proof should identify a common core or a controlled sequence of
cores, show that low-margin coordinates can be charged to an
\(\mathcal O(1/\eps)\)-volume region, and sum clearly inactive activations
using the outer geometric contraction.

\subsection{Priority 3: analyze the online FIFO advantage}

AESP-LOCAPPR is empirically stronger than AESP-LOCGD because it immediately
reuses coordinate updates.  The support-to-locality proof in
\cref{lem:support-to-locality} is direct for simultaneous LOCGD and for a
normalized priority rule, but not for arbitrary FIFO ordering.  The missing
queue lemma should quantify how much stale queue entries can inflate
\begin{equation}
    \frac{\overline{\vol}(S_t)}{\gamma_t}
    \label{eq:fifo-locality-object}
\end{equation}
relative to the support envelope.  A constant-factor comparison would preserve
the practical implementation while closing the theorem.

\subsection{Priority 4: improve the inner accuracy certificate}

The current \(1/\eps^2\) bound uses the mixed
\(\ell_\infty\)-by-\(\ell_1\) certificate
\cref{eq:mixed-norm-certificate}.  A local or incrementally maintainable
\(\ell_2\) certificate would have the desired square-root relation
\begin{equation}
    \norm{\grad h_t(z)}_2
    =\mathcal O(\sqrt{\phi_t}),
    \label{eq:l2-certificate-target}
\end{equation}
so a final \(\phi_t=\mathcal O(\eps^2)\) would require an
\(\mathcal O(\eps)\) gradient scale rather than
\(\mathcal O(\eps^2)\).  The challenge is to verify this without scanning the
whole graph.

\subsection{Priority 5: over-relaxation and accelerated tails}

There are three distinct questions:
\begin{enumerate}[leftmargin=2em]
\item Can the safeguard of \cref{def:safeguarded-sor} accept
      \(\omega_\star\) on most practical updates while preserving the mass
      theorem?
\item Can a block or sweep-level potential prove accelerated local behavior
      even though one-coordinate weighted \(\ell_1\) mass can increase?
\item Can a short verified LocCH block be used as an optional branch, with
      rollback to the proved \(\omega=1\) tail when its locality envelope or
      residual reduction fails?
\end{enumerate}
The current project already contains a rigorous long-spider obstruction for a
different two-stage method---coarse monotone forward push followed by fixed
optimal SOR.  That result should not be silently transferred to AESP, but it
warns that fixed-relaxation spectral intuition alone cannot prove a universal
local accelerated bound.

\subsection{Priority 6: arbitrary-warm-start locality for RPPR}

\Cref{thm:rppr-hybrid-convergence} proves that composite Catalyst followed by
full ISTA converges, but the COLT-style degree-work target remains open.  A
useful theorem should establish one of the following:
\begin{enumerate}[leftmargin=2em]
\item a certified support set \(\Omega\supseteq\supp(x_\rho^\star)\) of volume
      \(\mathcal O(1/\rho)\), followed by a zero restart on \(\Omega\);
\item a weighted proximal/KKT potential that decreases under local coordinate
      updates from an arbitrary signed or overshooting warm start; or
\item an over-regularized boundary-charging theorem whose cumulative overhead
      is compatible with \(\softO(1/(\rho\sqrt\alpha))\).
\end{enumerate}
The correct stopping object is the fixed-point residual \(R_\rho\), with the
solution certificate divided by \(\alpha\).  Any proof that instead reuses
zero-start support monotonicity without qualification is incomplete.

\subsection{Priority 7: benchmark implementation}

The TB-Science task can be built in parallel with the theoretical work once a
private hybrid reference reliably beats frozen AESP.  The implementation order
should be:
\begin{enumerate}[leftmargin=2em]
\item freeze the AESP baseline and graph-access instrumentation;
\item implement the proof-oriented hybrid and exact stopping certificate;
\item calibrate a hidden suite against all standalone baselines;
\item run frontier-agent pilots and set the speed threshold;
\item submit the concise proposal and disclose provenance;
\item only then expose the task repository and tests required by the Harbor
      format.
\end{enumerate}

\subsection{Precise claims that should not yet appear in the paper}

The following unconditional statements are not supported by the current
proof:
\begin{align}
    &\text{``AESP--LocSOR has }
    \softO(1/(\sqrt\alpha\eps))
    \text{ work on every graph,''}
    \label{eq:not-yet-claim}\\
    &\text{``Composite Catalyst--ISTA has }
    \softO(1/(\rho\sqrt\alpha))
    \text{ local work from every handoff.''}
    \label{eq:not-yet-rppr-claim}
\end{align}
The manuscript-ready unregularized statements are
\cref{thm:trajectory-total,thm:conditional-total,cor:black-box-rate}.  The
manuscript-ready RPPR statement is \cref{thm:rppr-hybrid-convergence}, followed
by an explicit description of the missing arbitrary-warm-start locality
lemma.  These distinctions are central to the scientific credibility of the
project.
